% 第1章 FLRW計量と膨張宇宙（仕様 01_textbook_spec.md §第1章）
\section{FLRW計量と膨張宇宙}
\label{ch:flrw}

\paragraph{この章で導くこと（入力→出力）}
一様等方性の仮定（入力）から FLRW 計量を書き下し、Einstein 方程式から
Friedmann 方程式（出力 S1.1）を導く。連続の式から各成分のスケーリング
（S1.2）、ヌル測地線から赤方偏移 $1+z=1/a$（S1.3）を得て、
$\Lambda$CDM 背景の $\mathcal H(x)$ と共形時間 $\eta(x)$（S1.4）を計算可能にする。
最終的に、$C_\ell$ の射影で用いる共動角径距離 $\chi(z_*)$ を数値評価する。

\subsection{宇宙原理と観測的根拠}
CMB の温度等方性は $\Delta T/T \sim 10^{-5}$ であり、これが本書全体の出発点
である「一様等方な背景＋微小摂動」という描像を正当化する。背景宇宙を本章で、
摂動を第\ref{ch:flrw}章以降で扱う。

\subsection{FLRW計量}
最大対称な3次元空間を空間切断にもつ時空の計量は、計量符号 $(-,+,+,+)$ で
\begin{equation}
  \mathrm ds^2 = -c^2\,\mathrm dt^2
  + a^2(t)\left[\frac{\mathrm dr^2}{1-k r^2}
  + r^2\,\mathrm d\Omega^2\right],
  \label{eq:flrw}
\end{equation}
と書ける。本書は主筋で平坦 $k=0$ を採り、曲率は第13章で復活させる。

\subsection{赤方偏移}
ヌル測地線に沿う光子の4元運動量は $p\propto 1/a$ を満たし（導出契約 S1.3）、
観測される波長の伸びは
\begin{equation}
  1+z = \frac{a_0}{a} = \frac{1}{a},\qquad a_0\equiv 1 .
\end{equation}

\subsection{Friedmann方程式（導出契約 S1.1）}
計量 \eqref{eq:flrw} の Christoffel 記号・Ricci テンソルを計算し（全成分は付録Aの表）、
Einstein 方程式 $G_{\mu\nu}=\frac{8\pi G}{c^4}T_{\mu\nu}$ の $00$ 成分から
\begin{equation}
  H^2 \equiv \left(\frac{\dot a}{a}\right)^2
  = \frac{8\pi G}{3}\rho - \frac{k c^2}{a^2},
  \label{eq:friedmann}
\end{equation}
空間成分との組み合わせから加速度方程式
$\ddot a/a = -\tfrac{4\pi G}{3}(\rho+3p/c^2)$ を得る。

\begin{intuition}
$H^2\propto\rho$ は「中身が多いほど速く膨張（または収縮）する」というニュートン
的描像と一致する。曲率項 $-kc^2/a^2$ は全エネルギー（運動＋ポテンシャル）の
符号に対応する。
\end{intuition}

\subsection{構成要素とスケーリング（導出契約 S1.2）}
共変保存則 $\nabla_\mu T^{\mu\nu}=0$ の時間成分は連続の式
$\dot\rho + 3H(\rho+p/c^2)=0$ を与え、状態方程式 $p=w\rho c^2$ に対し
$\rho\propto a^{-3(1+w)}$、すなわち
\begin{equation}
  \rho_r\propto a^{-4},\quad \rho_m\propto a^{-3},\quad \rho_\Lambda=\text{const}.
\end{equation}

\subsection{共形ハッブルと共形時間（導出契約 S1.4）}
$x\equiv\ln a$、$\mathcal H\equiv aH$ とすると、平坦宇宙で
\begin{equation}
  \mathcal H(x) = H_0\sqrt{\Omega_{m0}e^{-x}
  + \Omega_{r0}e^{-2x} + \Omega_{\Lambda 0}e^{2x}},
  \label{eq:Hp}
\end{equation}
共形時間（本書では長さ次元、$c$ 込み）は
$\eta(x)=\int_{-\infty}^{x}\frac{c\,\mathrm dx'}{\mathcal H(x')}$。
等密度は $a_{\rm eq}=\Omega_{r0}/\Omega_{m0}$、数値で
$z_{\rm eq}\approx\input{|"python3 -c \"import json;print(int(json.load(open('../figures/values.json'))['z_eq']))\""}$
（\texttt{values.json} 由来）。

\begin{numreality}
式 \eqref{eq:Hp} は \texttt{cmbcore.background.BackgroundCosmology.Hp} に
そのまま実装されており、$\eta(x),t(x)$ は ODE を後退/前進積分して
\texttt{CubicSpline} 化される。本章の全数値（$z_{\rm eq}$, $\eta_0$, $\chi(z_*)$）は
\texttt{scripts/make\_values.py} が生成する \texttt{values.json} から差し込まれる。
\end{numreality}

\subsection{距離}
共動距離 $\chi(x)=\eta_0-\eta(x)$。最終散乱面までの共動角径距離 $\chi(z_*)$ は
$C_\ell$ の角度射影 $\ell\sim k\chi$ を決める基本量であり、本書の fiducial
宇宙では $\chi(z_*)\approx 13.9$\,Gpc（第11章・第14章で使用）。

\paragraph{まとめ}
背景宇宙は $\mathcal H(x)$ ひとつで記述され、そこから $\eta,\chi,z_{\rm eq}$ が
すべて計算できる。次章ではこの背景上で熱史を組み立てる。

\paragraph{演習}
(1) 物質優勢で $\eta\propto a^{1/2}$ を示せ。
(2) $z_{\rm eq}$ の $\Omega_m h^2$ 依存を求めよ。
(3) NB1 で $\Omega_\Lambda$ を変え $\eta_0$ の変化を調べよ。
